\documentclass[a4paper,landscape,8pt]{extarticle}
\usepackage[landscape,margin=0.7cm,top=0.5cm,bottom=0.7cm]{geometry}
\usepackage[utf8]{inputenc}
\usepackage[T1]{fontenc}
\usepackage[ngerman]{babel}
\usepackage{helvet}
\renewcommand{\familydefault}{\sfdefault}
\usepackage{xcolor}
\usepackage{tikz}
\usetikzlibrary{positioning}
\usepackage{enumitem}
\usepackage{multicol}
\usepackage{array}
\usepackage{fancyhdr}
\usepackage{amssymb}

% Farben
\definecolor{photopink}{HTML}{E91E63}
\definecolor{albumblue}{HTML}{1565C0}
\definecolor{trashred}{HTML}{C62828}
\definecolor{lightgray}{HTML}{F5F5F5}
\definecolor{bordergray}{HTML}{BBBBBB}

\pagestyle{fancy}
\fancyhf{}
\renewcommand{\headrulewidth}{0pt}
\fancyfoot[C]{\footnotesize Seite \thepage{} von 3}

\setlength{\parindent}{0pt}
\setlength{\parskip}{2pt}

% Kompakte Listen
\setlist{nosep, leftmargin=1em, itemsep=0pt, topsep=1pt}

\begin{document}

% ============================================================================
% SEITE 1: INFORMATIONSBLATT
% ============================================================================

\noindent{\Large\bfseries Fotos \& Dateien: Ordnung in deinen Bildern} \hfill
{\footnotesize\textbf{Lernziele:} Fotos finden | Alben erstellen | Platz schaffen}

\vspace{1mm}
\hrule
\vspace{1mm}

% Drei-Spalten-Layout
\begin{minipage}[t]{0.32\linewidth}
\begin{center}
\colorbox{photopink!15}{%
\begin{minipage}{0.95\linewidth}
\vspace{2mm}
\begin{center}
{\Large\bfseries\textcolor{photopink}{1. Die Fotos-App}}\\[1pt]
{\footnotesize\itshape Alle Bilder an einem Ort}
\end{center}
\vspace{2mm}
\end{minipage}%
}
\end{center}

\vspace{1mm}

\textbf{Was ist das?}\\
Dein digitales \textbf{Fotoalbum} -- alle Bilder und Videos automatisch gesammelt.

\textbf{Die Bereiche:}
\begin{itemize}
\item \textbf{Mediathek:} Alle Fotos nach Datum
\item \textbf{Für dich:} Automatische Zusammenstellungen
\item \textbf{Alben:} Deine eigenen Sammlungen
\item \textbf{Suchen:} Fotos nach Inhalt finden
\end{itemize}

\textbf{Suche nutzen:}\\
Einfach Stichwort eingeben:
\begin{itemize}
\item Ort: ``Berlin'', ``Strand''
\item Zeit: ``2023'', ``Dezember''
\item Inhalt: ``Hund'', ``Geburtstag''
\end{itemize}

\textbf{Das iPhone erkennt automatisch}, was auf Fotos ist!

\end{minipage}%
\hfill%
\begin{minipage}[t]{0.32\linewidth}
\begin{center}
\colorbox{albumblue!10}{%
\begin{minipage}{0.95\linewidth}
\vspace{2mm}
\begin{center}
{\Large\bfseries\textcolor{albumblue}{2. Alben erstellen}}\\[1pt]
{\footnotesize\itshape Ordnung schaffen}
\end{center}
\vspace{2mm}
\end{minipage}%
}
\end{center}

\vspace{1mm}

\textbf{Warum Alben?}
\begin{itemize}
\item Fotos schneller finden
\item Themen zusammenfassen
\item Lieblingsfotos sammeln
\end{itemize}

\textbf{So erstellst du ein Album:}
\begin{enumerate}
\item ``Alben'' antippen (unten)
\item \textbf{+} oben links
\item ``Neues Album'' wählen
\item Namen eingeben (z.B. ``Enkel'')
\item Fotos auswählen
\item ``Hinzufügen''
\end{enumerate}

\textbf{Foto zu Album hinzufügen:}
\begin{enumerate}
\item Foto öffnen
\item Teilen-Symbol (Quadrat mit Pfeil)
\item ``Zum Album hinzufügen''
\item Album auswählen
\end{enumerate}

\textbf{Album-Ideen:}
\begin{itemize}
\item Familie / Enkel
\item Urlaub 2024
\item Garten
\item Rezepte (abfotografiert)
\end{itemize}

\end{minipage}%
\hfill%
\begin{minipage}[t]{0.32\linewidth}
\begin{center}
\colorbox{trashred!10}{%
\begin{minipage}{0.95\linewidth}
\vspace{2mm}
\begin{center}
{\Large\bfseries\textcolor{trashred}{3. Löschen \& Platz}}\\[1pt]
{\footnotesize\itshape Speicher voll?}
\end{center}
\vspace{2mm}
\end{minipage}%
}
\end{center}

\vspace{1mm}

\textbf{Foto löschen:}
\begin{enumerate}
\item Foto(s) auswählen
\item Papierkorb-Symbol unten
\item ``Foto löschen'' bestätigen
\end{enumerate}

\textbf{Wichtig: ``Zuletzt gelöscht''}
\begin{itemize}
\item Gelöschte Fotos sind \textbf{30 Tage} noch da!
\item Unter ``Alben'' $\rightarrow$ ``Zuletzt gelöscht''
\item Von dort: \textbf{Wiederherstellen} möglich
\item Oder: Endgültig löschen
\end{itemize}

\vspace{2mm}
\fcolorbox{trashred}{trashred!5}{%
\begin{minipage}{0.92\linewidth}
\footnotesize
\textbf{Tipp:} ``Zuletzt gelöscht'' leeren spart Speicher!
\end{minipage}%
}

\vspace{2mm}

\textbf{Speicher-Tipps:}
\begin{itemize}
\item Duplikate löschen (Alben $\rightarrow$ Duplikate)
\item Große Videos prüfen
\item Unscharfe Fotos löschen
\item iCloud nutzen (siehe Einheit 018)
\end{itemize}

\end{minipage}

\vspace{1mm}

% Zusammenfassung
\begin{center}
\fcolorbox{bordergray}{lightgray}{%
\begin{minipage}{0.98\linewidth}
\centering
\vspace{1.5mm}
\textbf{Zusammenfassung:}
\textcolor{photopink}{\textbf{Suche}} = Ort, Zeit, Inhalt eingeben ~$|$~
\textcolor{albumblue}{\textbf{Album}} = Fotos thematisch sammeln ~$|$~
\textcolor{trashred}{\textbf{Zuletzt gelöscht}} = 30 Tage Rettungschance
\vspace{1.5mm}
\end{minipage}%
}
\end{center}

\newpage

% ============================================================================
% SEITE 2: ÜBUNGEN
% ============================================================================

\begin{center}
{\LARGE\bfseries Übungen: Fotos organisieren}\\[3pt]
{\normalsize Schau dir Seite 1 an und beantworte die Fragen}
\end{center}

\vspace{3mm}
\hrule
\vspace{4mm}

\begin{multicols}{2}

\textbf{\large Teil A: Grundlagen verstehen}

Richtig oder Falsch? Kreuze an.

\vspace{2mm}

\renewcommand{\arraystretch}{1.6}
\begin{tabular}{@{}p{8cm}cc@{}}
& R & F \\
1. Die Fotos-App sammelt alle Bilder automatisch. & $\square$ & $\square$ \\
2. Man kann nach ``Strand'' oder ``Hund'' suchen. & $\square$ & $\square$ \\
3. Gelöschte Fotos sind sofort endgültig weg. & $\square$ & $\square$ \\
4. Alben helfen, Fotos zu organisieren. & $\square$ & $\square$ \\
5. ``Zuletzt gelöscht'' speichert Fotos 30 Tage. & $\square$ & $\square$ \\
\end{tabular}
\renewcommand{\arraystretch}{1}

\vspace{6mm}

\textbf{\large Teil B: Suche verstehen}

Welches Suchwort würdest du eingeben? Schreibe es auf.

\vspace{2mm}

\renewcommand{\arraystretch}{1.8}
\begin{tabular}{@{}p{6cm}p{5cm}@{}}
1. Fotos vom Urlaub am Meer: & \dotfill \\
2. Fotos von letztem Weihnachten: & \dotfill \\
3. Fotos mit deinem Hund: & \dotfill \\
4. Fotos aus Berlin: & \dotfill \\
\end{tabular}
\renewcommand{\arraystretch}{1}

\vspace{6mm}

\textbf{\large Teil C: Album erstellen}

Nummeriere die Schritte (1--5):

\vspace{2mm}

\renewcommand{\arraystretch}{1.8}
\begin{tabular}{@{}cp{7cm}@{}}
$\square$ & ``Neues Album'' wählen \\
$\square$ & Namen eingeben (z.B. ``Familie'') \\
$\square$ & Fotos auswählen und hinzufügen \\
$\square$ & ``Alben'' unten antippen \\
$\square$ & + Symbol antippen \\
\end{tabular}
\renewcommand{\arraystretch}{1}

\columnbreak

\textbf{\large Teil D: Praktisch üben}

Führe diese Aufgaben auf deinem Gerät durch:

\vspace{3mm}

\textbf{Aufgabe 1: Suchen}\\
Suche nach einem Ort, an dem du Fotos gemacht hast.

\vspace{2mm}
Was hast du eingegeben? \dotfill

\vspace{2mm}
Wie viele Fotos wurden gefunden? \dotfill

\vspace{5mm}

\textbf{Aufgabe 2: Album erstellen}\\
Erstelle ein Album mit dem Namen ``Lieblingsfotos'' und füge mindestens 3 Fotos hinzu.

\vspace{2mm}
$\square$ Erledigt

\vspace{5mm}

\textbf{Aufgabe 3: Zuletzt gelöscht prüfen}\\
Schau nach, ob Fotos in ``Zuletzt gelöscht'' sind.

\vspace{2mm}
Wie viele sind dort? \dotfill

\vspace{5mm}

\textbf{Bonus: Speicher prüfen}\\
Gehe zu Einstellungen $\rightarrow$ Allgemein $\rightarrow$ iPhone-Speicher

\vspace{2mm}
Wie viel Speicher nutzen Fotos? \dotfill

\end{multicols}

\newpage

% ============================================================================
% SEITE 3: CHECKLISTE
% ============================================================================

\begin{center}
{\LARGE\bfseries Checkliste: Das habe ich verstanden}\\[3pt]
{\normalsize Geh die Liste durch und hake ab, was du sicher weißt}
\end{center}

\vspace{3mm}
\hrule
\vspace{6mm}

% Drei Boxen nebeneinander
\begin{minipage}[t]{0.31\linewidth}
\begin{center}
\fcolorbox{photopink}{photopink!10}{%
\begin{minipage}{0.92\linewidth}
\vspace{3mm}
\begin{center}
{\large\bfseries\textcolor{photopink}{Fotos finden}}
\end{center}
\vspace{2mm}
\begin{itemize}[label=$\square$, itemsep=5pt]
\item Ich kann die \textbf{Fotos-App} öffnen
\item Ich kenne die \textbf{Bereiche} (Mediathek, Alben...)
\item Ich kann Fotos \textbf{suchen} (Ort, Zeit, Inhalt)
\item Ich weiß, dass das iPhone Inhalte \textbf{erkennt}
\end{itemize}
\vspace{3mm}
\end{minipage}%
}
\end{center}
\end{minipage}%
\hfill%
\begin{minipage}[t]{0.31\linewidth}
\begin{center}
\fcolorbox{albumblue}{albumblue!8}{%
\begin{minipage}{0.92\linewidth}
\vspace{3mm}
\begin{center}
{\large\bfseries\textcolor{albumblue}{Alben nutzen}}
\end{center}
\vspace{2mm}
\begin{itemize}[label=$\square$, itemsep=5pt]
\item Ich kann ein \textbf{Album erstellen}
\item Ich kann Fotos \textbf{zu Alben hinzufügen}
\item Ich habe Ideen für \textbf{nützliche Alben}
\item Alben helfen mir, \textbf{Ordnung} zu halten
\end{itemize}
\vspace{3mm}
\end{minipage}%
}
\end{center}
\end{minipage}%
\hfill%
\begin{minipage}[t]{0.31\linewidth}
\begin{center}
\fcolorbox{trashred}{trashred!8}{%
\begin{minipage}{0.92\linewidth}
\vspace{3mm}
\begin{center}
{\large\bfseries\textcolor{trashred}{Löschen \& Speicher}}
\end{center}
\vspace{2mm}
\begin{itemize}[label=$\square$, itemsep=5pt]
\item Ich kann Fotos \textbf{löschen}
\item Ich kenne \textbf{``Zuletzt gelöscht''}
\item Ich kann Fotos \textbf{wiederherstellen}
\item Ich weiß, wie ich \textbf{Speicher spare}
\end{itemize}
\vspace{3mm}
\end{minipage}%
}
\end{center}
\end{minipage}

\vspace{8mm}

% Gesamtverständnis
\begin{center}
\fcolorbox{black}{lightgray}{%
\begin{minipage}{0.95\linewidth}
\vspace{4mm}
\begin{center}
{\large\bfseries Gesamtverständnis}
\end{center}
\vspace{3mm}
\begin{multicols}{2}
\begin{itemize}[label=$\square$, itemsep=6pt]
\item Ich kann \textbf{Fotos finden} durch die Suche
\item Ich kann \textbf{Alben} für meine Themen erstellen
\item Ich weiß, wie ich Fotos \textbf{sicher lösche}
\item Ich kann gelöschte Fotos \textbf{wiederherstellen}
\item Ich weiß, wie ich \textbf{Speicherplatz} spare
\end{itemize}
\end{multicols}
\vspace{3mm}
\end{minipage}%
}
\end{center}

\vspace{8mm}

% Notfall-Notizen
\begin{center}
\fcolorbox{bordergray}{white}{%
\begin{minipage}{0.95\linewidth}
\vspace{4mm}
\begin{center}
{\large\bfseries Meine Foto-Organisation}\\[3pt]
{\footnotesize (Welche Alben möchtest du erstellen?)}
\end{center}
\vspace{4mm}
\renewcommand{\arraystretch}{2}
\begin{tabular}{@{}p{0.45\linewidth}p{0.45\linewidth}@{}}
\textbf{Geplante Alben:} & \textbf{Mein Speicherstand:} \\
1. \dotfill & Gesamt: \dotfill GB \\[2mm]
2. \dotfill & Fotos nutzen: \dotfill GB \\[2mm]
3. \dotfill & Noch frei: \dotfill GB \\[2mm]
4. \dotfill & \\
\end{tabular}
\renewcommand{\arraystretch}{1}
\vspace{4mm}
\end{minipage}%
}
\end{center}

\vfill

\begin{center}
\footnotesize\textit{Stand: \today{} -- Regelmäßig aufräumen: Duplikate und unscharfe Fotos löschen spart Speicher!}
\end{center}

\end{document}
